% Chapter Template

\chapter{Literature Review}\singlespacing % Main chapter title

\label{Chapter2} % Change X to a consecutive number; for referencing this chapter elsewhere, use \ref{ChapterX}

\lhead{Chapter 2. \emph{Literature Review}} % 

\section{Neurons in Multilingual Models}
Understanding the representation of language-specific and cross-lingual knowledge in multilingual language models (ML-LMs) is crucial for improving their generalization capabilities across diverse languages. Recent studies, such as those probing mBERT and XLM-R using the mLAMA dataset [\cite{zhao2024tracing}], have revealed that factual knowledge in ML-LMs is encoded in three distinct patterns: language-independent, cross-lingual shared, and cross-lingual transferred representations. Language-independent neurons encode facts uniquely for each language, whereas cross-lingual shared neurons represent the same facts using identical neural patterns across languages. Cross-lingual transferred neurons leverage knowledge from high-resource languages to infer facts in low-resource languages. However, limitations in current models hinder consistent fact representation and effective cross-lingual transfer, underscoring the need for targeted interventions and better probing techniques to enhance their multilingual factual capabilities.

The study by \cite{kojima2024} explored language-specific neurons in decoder-based pre-trained language models (PLMs) such as XGLM, BLOOM, and Llama2. These neurons, which are uniquely activated for specific languages, were found predominantly in the initial and final layers of the models. The researchers demonstrated that manipulating these neurons during inference could significantly influence the language probability of generated text. The findings highlight the intricate role of language-specific neurons in cross-lingual text generation, suggesting that decoder-based PLMs exhibit a layered structure of multilingual capabilities, where early and late layers handle language-specific processing while middle layers remain language-agnostic. This work underscores the potential of targeted neuron interventions to enhance multilingual performance in PLMs.

Task-specific neurons \cite{leng2024} in large language models (LLMs) have been empirically shown to play a critical role in multi-task learning and generalization. Using gradient attribution, these neurons are identified as those with the highest relevance to specific tasks, and their deactivation or fine-tuning significantly impacts task-specific performance. Importantly, overlaps in task-specific neurons across tasks are strongly correlated with generalization capabilities, providing evidence of shared knowledge representations in LLMs. Additionally, methods like Neuron-level Continuous Fine-Tuning (NCFT) leverage these neurons to mitigate catastrophic forgetting in continuous learning scenarios by isolating and dynamically weighting task-specific parameters during training and inference, as demonstrated in recent studies.

Recent advancements in understanding multilingual capabilities in LLMs emphasize the pivotal role of language-specific neurons. \cite{tang2024languagespecific} propose the novel Language Activation Probability Entropy (LAPE) metric to identify neurons that are preferentially activated by specific languages, enabling insights into the organization of multilingual capabilities within models like LLaMA-2 and BLOOM. Their findings reveal that these neurons are predominantly located in the top and bottom layers, facilitating both the transformation of language-specific inputs into shared semantic representations and the generation of outputs in respective languages. Furthermore, selective activation or deactivation of these neurons demonstrates the ability to steer model responses in a particular language, showcasing the potential for targeted interventions to enhance cross-lingual performance.

Studies on LLMs have revealed the presence of language-specific neurons, which play a crucial role in handling multilingual tasks. These neurons exhibit distinct activation patterns when processing inputs in specific languages, as identified through methods such as Parallel Language-Specific Neuron Detection (PLND). The proposed Multilingual Workflow (MWork) hypothesizes that LLMs process queries in three stages: understanding the query by converting it into English, reasoning in English while incorporating multilingual knowledge through feed-forward structures, and finally generating responses in the original query language. Experiments demonstrate that fine-tuning these neurons can significantly enhance multilingual capabilities across both high-resource and low-resource languages without negatively impacting performance in other languages. These findings underscore the importance of language-specific neurons in improving multilingual generalization and task performance \cite{zhao2024large}.

Advancements in analyzing LLMs have identified the presence of style-specific neurons, which extend to multilingual capabilities. These neurons, as highlighted by \cite{lai2024stylespecific}, are crucial for tasks such as text style transfer (TST) and language adaptation. Style-specific neurons are identified by their selective activations for particular attributes, such as language or stylistic features, and can be manipulated to steer LLM outputs. For instance, deactivating source-specific neurons enhances the generation of target-style words but can impact fluency. Methods like contrastive decoding, adapted for style-specific tasks, mitigate these effects by refining token probabilities across layers. Such neuron-level analysis not only demonstrates the modularity within LLMs but also underscores their adaptability for multilingual and stylistic transformations.

Recent research has revealed the critical role of neurons in the factual knowledge representation and retrieval within multilingual LLMs. For example, \cite{meng2022locating} identified specific mid-layer feedforward networks in GPT models that localize factual associations, which can be directly edited without disrupting unrelated facts. Their work proposed Rank-One Model Editing (ROME), a method to modify the memory of such neurons effectively, demonstrating high generalization and specificity for knowledge updates. This study underscores the importance of mid-layer computations for cross-lingual tasks, enabling more interpretable and controllable model adaptations. 

\section{Mechanistic Interpretability and Multilinguality}

One notable contribution in the field of mechanistic interpretability is the introduction of Softmax Linear Units (SoLU), an architectural modification designed to enhance neuron interpretability. By replacing traditional activation functions with SoLU, researchers observed a substantial increase in neurons that correspond to human-understandable concepts, phrases, or categories. This enhancement facilitates a clearer understanding of how models process information, including multilingual data, by making the internal representations more transparent and interpretable [\cite{elhage2022solu}].

Building on this, Anthropic employed dictionary learning techniques to decompose language models into understandable components. By identifying patterns of neuron activations, termed "features," researchers could map these to specific concepts within the model. This approach was successfully scaled to larger models like Claude 3.0 Sonnet, revealing features corresponding to a vast range of entities and concepts across multiple languages. These findings underscore the model's capacity to represent multilingual information through distinct, interpretable features, enhancing our understanding of how LLMs manage multilinguality [\cite{bricken2023monosemanticity}].

Additionally, Anthropic's exploration into the phenomenon of superposition—where models represent more features than they have dimensions—has implications for multilingual processing. Understanding how models compress and manage multiple linguistic features within limited dimensions can inform strategies to improve multilingual performance and interpretability. By dissecting these mechanisms, researchers can develop more efficient models capable of handling multiple languages without compromising interpretability [\cite{elhage2022superposition}].










